\subsection{Choosing the right collection (E3)}
\label{sec:routing-audit}
We first asked whether the description on a catalogue card helps the
selector find the right collection. A0 supplied the name and a description
derived from the collection's contents; A3 was intended to supply the name
alone. With descriptions, the scripted selector chose the labelled
collection on all 72 questions, compared with 37 under A3 in the frozen run
(Figure~\ref{fig:routing}). The direction persisted on the two reruns.
Even with the collection correctly identified, A0 returned the target
string on only 37 questions. Searching in the right place does not ensure
that the right passage reaches the reader.

\begin{figure}[!htb]
\centering
\includegraphics[width=\columnwidth]{fig_opt_route}
\caption{Correct collection selections for the two route-card variants.
All three runs use the same 72 questions. A3 is nominally name-only.}
\label{fig:routing}
\end{figure}

\subsection{Getting the target text into the evidence pack (E4)}
\label{sec:ladder-opt}
On the same questions, we compared the six ways of assembling evidence in
Table~\ref{tab:ablation-configs}. In the frozen run, A0h retrieved the
annotated string on 52 of 72 questions, up from 37 with A0. It recovered 15
of A0's misses without losing any of its hits (exact McNemar
$p=6.1\times10^{-5}$). Both same-question reruns preserved the gain.
A0h also led the unrouted vector union A1b, but the difference from the
unrouted per-document configuration PD was not significant
(Figure~\ref{fig:content}).

A0h changes what the system retrieves and what it puts first in the pack.
The comparison with A0f is therefore not a test of window size alone.
Similarly, A2 removes document reads as well as the gate, so it does not
isolate the gate's effect. The trace counted 309~kB of text per question
for A0h, against 31.4~kB for A0 and 24.3~MB for the mechanical read-all scan,
which found every target. For A0h the count sums candidate text fetched for
the gate, the gate's output and selected chunk windows. Text can be counted
more than once. The final A0h pack has a separate 60{,}000-character cap;
these byte totals are neither final pack sizes nor generator input sizes.

\begin{figure*}[!tb]
\centering
\includegraphics[width=\textwidth]{fig_opt_content}
\caption{Target-string containment on 72 in-domain questions. (a) Circles
and Wilson 95\% intervals show the frozen run; open triangles show the
same-question follow-up. The dashed line marks frozen A0. The columns give
hit counts out of 72; a dash indicates that the read-all follow-up was
incomplete. (b) For each frozen-run contrast, the symbols count questions
on which only the first or only the second configuration hit. McNemar
$p$ values are exact, two-sided and unadjusted.}
\label{fig:content}
\end{figure*}

\subsection{Storing and finding operational lessons (E5)}
\label{sec:e3}
HMRS stores operational lessons as searchable documents and orders them by
credibility tier. The extraction check produced draft lessons, and one
stored lesson ranked first in a case-level search. No generated response
was retained for that case. The approval check found an empty draft pool,
and the heuristic dry run did not meet its quality criteria. These checks
show that the system can store and retrieve lessons; they do not establish
reliable automatic approval, feedback or credibility decay over time.
